% !TeX root = ../../Book4.tex
% 纲要标题：### 9.3 Tor
% 纲要标签：b4:sec:0803
% 纲要位置：计划纲要.md，第 3414 行。

\section{Tor}
\label{b4:sec:0803}

张量积的两个模必须取相反的侧别，任一变量的投射消解都应得到同一组 Tor。总复形的比较同时解释平衡性和短正合列中的连接映射。


% 纲要内容（仅作写作计划，尚非教材正文）：
%
% * 张量积的左导出函子
% * 平衡性
% * Tor 长正合列及具体计算
%

\subsection{张量积的左导出函子}
\label{b4:subsec:0803:01}

\begin{definition}\label{b4:def:tor}
设 \(R\) 是含幺环，\(M\) 是幺性右 \(R\) 模，\(N\) 是幺性左 \(R\) 模。对 \(M\) 取右模投射消解 \(P_\bullet\to M\)，定义交换群
\[
\operatorname{Tor}_n^R(M,N)\coloneqq H_n(P_\bullet\otimes_R N),\qquad n\ge0.
\]
这就是右正合函子 \(-\otimes_R N\) 的左导出函子。
\end{definition}
\symidx{\(\operatorname{Tor}_n^R(M,N)\)：张量积函子的左导出函子值}

投射消解在零次给出 \(P_1\to P_0\to M\to0\)。张量的右正合性使
\[
P_1\otimes_RN\longrightarrow P_0\otimes_RN
\longrightarrow M\otimes_RN\longrightarrow0
\]
正合，因此
\[
\operatorname{Tor}_0^R(M,N)
=\operatorname{coker}(P_1\otimes_RN\to P_0\otimes_RN)
\cong M\otimes_RN.
\]
这一零次识别只用右正合性，不要求 \(N\) 平坦。

\ProseParagraphBreak
左右侧别不可交换省略。若 \(P_\bullet\to M\) 和 \(Q_\bullet\to N\) 分别为右模、左模投射消解，下一小节的平衡性给出
\[
H_n(P_\bullet\otimes_RN)\cong H_n(M\otimes_RQ_\bullet).
\]
它没有交换 \(M,N\)，只改变在哪个变量取消解。只有在交换环上，左右模有自然识别，才能进一步比较 \(\operatorname{Tor}_n^R(M,N)\) 与 \(\operatorname{Tor}_n^R(N,M)\)，得到对称性。

\subsection{平衡性}
\label{b4:subsec:0803:02}

\begin{theorem}[Tor 的平衡性]\label{b4:thm:tor-balanced}
设 \(R\) 是含幺环，\(M\) 是幺性右 \(R\) 模，\(N\) 是幺性左 \(R\) 模。若 \(Q_\bullet\to N\) 是左 \(R\) 模投射消解，则对每个整数 \(n\ge0\)，有对两个变量自然的同构
\[
\operatorname{Tor}_n^R(M,N)\cong H_n(M\otimes_RQ_\bullet).
\]
因而可以在任意一侧取投射消解计算 Tor。
\end{theorem}
\begin{proof}
与 Ext 的平衡性一样，取两份消解组成双复形，再比较其两条边缘。这里使行、列正合的性质是投射模的平坦性。取右模投射消解 \(P_\bullet\to M\)，令 \(D_{p,q}=P_p\otimes_RQ_q\)，\(p,q\ge0\)，以
\[
\mathrm{d}_h=\mathrm{d}_P\otimes1,\qquad \mathrm{d}_v=(-1)^p1\otimes \mathrm{d}_Q
\]
为微分。它们反交换，所以总链微分平方为零。第二卷定理7.3.4证明了投射左模平坦。投射右 \(R\) 模正是投射左 \(R^{\mathrm{op}}\) 模，将该定理用于 \(R^{\mathrm{op}}\)，并经 \(X\otimes_{R^{\mathrm{op}}}P\cong P\otimes_RX\) 识别相应的张量函子，便知 \(P\otimes_R-\) 正合，即投射右模平坦。固定 \(p\) 时，\(P_p\) 是投射右 \(R\) 模，因而 \(P_p\otimes_R-\) 在左模范畴上正合；竖列只在零次有同调 \(P_p\otimes_RN\)。固定 \(q\) 时，\(Q_q\) 是投射左 \(R\) 模，因而 \(-\otimes_RQ_q\) 在右模范畴上正合；横行只在零次有同调 \(M\otimes_RQ_q\)。

记两份消解的增广为 \(\varepsilon_P:P_0\to M\)、\(\varepsilon_Q:Q_0\to N\)。第一条总化映射在 \(q=0\) 上取 \(1\otimes\varepsilon_Q\)，在 \(q>0\) 上取零；第二条在 \(p=0\) 上取 \(\varepsilon_P\otimes1\)，在 \(p>0\) 上取零。由增广零化对应的一次微分，它们与总链微分交换，给出
\[
\operatorname{Tot}D\longrightarrow P_\bullet\otimes_RN,
\qquad
\operatorname{Tot}D\longrightarrow M\otimes_RQ_\bullet.
\]
把两条边缘分别放在 \(q=0\) 和 \(p=0\)，刚才的纵列、横行同调计算说明这两条映射分别逐列、逐行拟同构。对于 \(n\ge0\)，总次数对角线 \(p+q=n\) 只有 \((0,n),(1,n-1),\ldots,(n,0)\) 共 \(n+1\) 个位置；\(n<0\) 时没有位置。因此\cref{b4:cor:double-complex-comparison}的链版本适用，两条总化映射均为拟同构。取同调并组合两同构即得结论。比较映射与增广交换，比较选择同伦，故所得同构对两变量自然。

还需要使两种计算给出相同的连接态射。对于第一变量的短正合列，取逐次分裂的右模马蹄消解，并固定第二变量的左模投射消解。三个总复形及边缘 \(P_\bullet\otimes_RN\) 的短正合列使用马蹄的逐次分裂性；另一边缘 \(M\otimes_RQ_\bullet\) 的短正合列则使用各 \(Q_q\) 平坦。对于第二变量的短正合列，改用左模马蹄消解并固定 \(P_\bullet\)，两条边缘的正合性分别来自各 \(P_p\) 平坦和马蹄的逐次分裂性。两种情况下，增广映射都给出总复形短正合列到边缘短正合列的交换图。由同调长正合列的自然性，平衡同构便与两变量的连接态射相容。
\end{proof}

若 \(R\) 交换，对左模定义右作用 \(n\cdot r\coloneqq rn\)，对右模定义左作用 \(r\cdot m\coloneqq mr\)，便可自然识别左右模；同样识别消解中的各项。这时交换两因子的映射才有如下张量积类型，并且满足平衡关系：
\[
\tau_{p,q}:P_p\otimes_RQ_q\longrightarrow Q_q\otimes_RP_p,
\qquad x\otimes y\longmapsto(-1)^{pq}y\otimes x.
\]
这些 \(\tau_{p,q}\) 合为保持总次数的映射 \(\tau\)。在交换后以 \(Q_q\) 为第一因子的总复形中，微分为 \(\mathrm{d}'=\mathrm{d}_Q\otimes1+(-1)^q1\otimes\mathrm{d}_P\)。对 \(x\in P_p\)、\(y\in Q_q\)，有
\[
\begin{aligned}
\mathrm{d}'\tau(x\otimes y)
&=(-1)^{pq}\cdot\mathrm{d}_Qy\otimes x
 +(-1)^{pq+q}y\otimes\mathrm{d}_Px,\\
\tau\mathrm{d}(x\otimes y)
&=(-1)^{(p-1)q}y\otimes\mathrm{d}_Px
 +(-1)^{p+p(q-1)}\cdot\mathrm{d}_Qy\otimes x.
\end{aligned}
\]
其中 \(p+p(q-1)=pq\)，而 \((p-1)q\) 与 \(pq+q\) 相差 \(2q\)，故两个分量的系数分别相同，得到 \(\mathrm{d}'\tau=\tau\mathrm{d}\)。因此 \(\tau\) 给出总复形同构，取同调并用平衡性，得到对两变量自然的同构 \(\operatorname{Tor}_n^R(M,N)\cong\operatorname{Tor}_n^R(N,M)\)。一般非交换环上，仅凭给定的左右模结构，不能作上述自然识别或定义这样的交换映射；平衡性仍只表示可在任一变量取消解。

\subsection{Tor 长正合列及具体计算}
\label{b4:subsec:0803:03}

\begin{theorem}\label{b4:thm:tor-long-exact}
设 \(R\) 是含幺环，\(M\) 是幺性右 \(R\) 模，\(N\) 是幺性左 \(R\) 模。幺性右模短正合列 \(0\to M'\to M\to M''\to0\) 给出
\[
\begin{aligned}
\cdots&\longrightarrow\operatorname{Tor}_1^R(M,N)
\longrightarrow\operatorname{Tor}_1^R(M'',N)
\longrightarrow M'\otimes_RN\\
&\longrightarrow M\otimes_RN\longrightarrow M''\otimes_RN\longrightarrow0,
\end{aligned}
\]
其一般中间段为
\[
\begin{aligned}
\cdots\longrightarrow\operatorname{Tor}_n^R(M',N)
&\longrightarrow\operatorname{Tor}_n^R(M,N)
\longrightarrow\operatorname{Tor}_n^R(M'',N)\\
&\xrightarrow{\delta}\operatorname{Tor}_{n-1}^R(M',N)
\longrightarrow\cdots\qquad(n\ge1).
\end{aligned}
\]
幺性左 \(R\) 模短正合列 \(0\to N'\to N\to N''\to0\) 同样给出自然长正合列，其一般中间段为
\[
\begin{aligned}
\cdots\longrightarrow\operatorname{Tor}_n^R(M,N')
&\longrightarrow\operatorname{Tor}_n^R(M,N)
\longrightarrow\operatorname{Tor}_n^R(M,N'')\\
&\xrightarrow{\delta}\operatorname{Tor}_{n-1}^R(M,N')
\longrightarrow\cdots\qquad(n\ge1).
\end{aligned}
\]
\end{theorem}
\begin{proof}
在第一变量取马蹄消解，得到逐次分裂的短正合列 \(0\to P_n'\to P_n\to P_n''\to0\)。加性函子 \(-\otimes_RN\) 保持分裂短正合列，故张量后仍为复形短正合列，可以取同调长正合列。这里保证左端单射的是分裂性；张量仅有右正合性并不足以得到它。第二变量可用平衡性改在该侧取马蹄消解，或固定右模投射消解并用各 \(P_n\) 平坦。由\cref{b4:thm:tor-balanced}证明中的连接态射相容性，两种构造对应同一长正合列；零次识别给出所列张量尾端。
\end{proof}

\begin{example}\label{b4:exm:tor-cyclic}
设 \(R\) 是主理想整环，\(0\ne a\in R\)，\(N\) 是幺性 \(R\) 模。将 \(R/(a)\) 的两项消解与 \(N\) 张量，得到链次数一、零上的 \(N\xrightarrow{\times a}N\)，故零次同调是乘 \(a\) 的余核，一次同调是其核：
\[
\operatorname{Tor}_0^R(R/(a),N)=N/aN,
\qquad \operatorname{Tor}_1^R(R/(a),N)=\{\,n\in N\mid an=0\,\},
\]
且高次为零。对 \(R=\mathbb Z\)、\(N=\mathbb Z/n\mathbb Z\)、\(a=m\)，其中 \(m,n>0\)，一次 Tor 同构于 \(\mathbb Z/\gcd(m,n)\mathbb Z\)。与\cref{b4:exm:ext-cyclic}比较，对同一个乘法映射 \(a:N\to N\)，\(\operatorname{Ext}_R^1(R/(a),N)\) 取其余核 \(N/aN\)，\(\operatorname{Tor}_1^R(R/(a),N)\) 则取其核。有限循环群算例给出同构的交换群，并不使两者成为同一构造。
\end{example}

设 \(k\) 为域，令 \(R=k[x,y]\)，并把 \(k\) 看作商模 \(R/(x,y)\)。\cref{b4:prop:koszul-resolution}给出的两变量消解中，\(P_0=R\) 记录一个生成元，\(P_1=R^2\) 的微分为 \((a,b)\mapsto xa+yb\)，记录两个关系 \(x,y\)，而 \(P_2=R\) 的微分为 \(c\mapsto(-yc,xc)\)，记录这两个关系之间的关系 \(x(-y)+yx=0\)。与 \(k\) 张量后，\(x,y\) 都作用为零，所以微分全为零，各项 \(k,k^2,k\) 直接成为同调：
\[
\operatorname{Tor}_0^R(k,k)\cong k,\qquad
\operatorname{Tor}_1^R(k,k)\cong k^2,\qquad
\operatorname{Tor}_2^R(k,k)\cong k,
\]
更高次为零。这里一次和二次 Tor 的维数分别记录了这个消解中的两个关系以及它们之间的一条关系。
